% Part: incompleteness
% Chapter: representability-in-q
% Section: comp-representable

\documentclass[../../../include/open-logic-section]{subfiles}

\begin{document}

\olfileid{inc}{req}{crq}
\olsection{গণনসাধ্য অপেক্ষকগুলি $\Th{Q}$-তে উপস্থাপনযোগ্য}

\begin{thm}
প্রত্যেক গণনসাধ্য অপেক্ষক $\Th{Q}$-তে উপস্থাপনযোগ্য।
\end{thm}

\begin{proof}
নির্দিষ্টতার জন্য এবং চার্চ--টুরিং অভিসন্দর্ভ ব্যবহার করে বলি, কোনো অপেক্ষক
গণনসাধ্য যদি এবং কেবল যদি সেটি সাধারণ পুনরাবৃত্ত। সাধারণ পুনরাবৃত্ত
অপেক্ষকগুলি হলো সেগুলি, যাদের শূন্য অপেক্ষক~$\Zero$, উত্তরসূরি
অপেক্ষক~$\Succ$ ও অভিক্ষেপ অপেক্ষক~$\Proj{n}{i}$ থেকে মিশ্রণ, আদিম
পুনরাবৃত্তি ও নিয়মিত ন্যূনীকরণের সাহায্যে সংজ্ঞায়িত করা যায়।
\olref[pri]{lem:prim-rec} অনুসারে, $f$ ও~$g$ থেকে আদিম পুনরাবৃত্তি দিয়ে
সংজ্ঞায়িত যেকোনো অপেক্ষক~$h$-কে $f$, $g$, $\Zero$, $\Succ$,
$\Proj{n}{i}$, $\Add$, $\Mult$, $\Char{=}$ থেকে মিশ্রণ ও নিয়মিত
ন্যূনীকরণের সাহায্যেও সংজ্ঞায়িত করা যায়। ফলত, কোনো অপেক্ষক সাধারণ
পুনরাবৃত্ত যদি এবং কেবল যদি $\Zero$, $\Succ$, $\Proj{n}{i}$, $\Add$,
$\Mult$, $\Char{=}$ থেকে মিশ্রণ ও নিয়মিত ন্যূনীকরণের সাহায্যে তাকে
সংজ্ঞায়িত করা যায়।

আরও দেখানো হয়েছে, আলোচ্য মৌলিক অপেক্ষকগুলি $\Th{Q}$-তে উপস্থাপনযোগ্য
(\cref{inc:req:bre:prop:rep-zero,inc:req:bre:prop:rep-succ,inc:req:bre:prop:rep-proj,inc:req:bre:prop:rep-id,inc:req:bre:prop:rep-add,inc:req:bre:prop:rep-mult}),
এবং উপস্থাপনযোগ্য অপেক্ষক থেকে মিশ্রণ বা নিয়মিত ন্যূনীকরণের সাহায্যে
সংজ্ঞায়িত যেকোনো অপেক্ষকও
(\olref[cmp]{prop:rep-composition},
\olref[min]{prop:rep-minimization}) উপস্থাপনযোগ্য। তাই প্রত্যেক সাধারণ
পুনরাবৃত্ত অপেক্ষক $\Th{Q}$-তে উপস্থাপনযোগ্য।
\end{proof}

\begin{explain}
আমরা দেখিয়েছি, গণনসাধ্য অপেক্ষকের সেটকে $\Th{Q}$-তে উপস্থাপনযোগ্য অপেক্ষকের
সেট হিসেবে চিহ্নিত করা যায়। বস্তুত প্রমাণটি আরও সাধারণ। উপস্থাপনযোগ্যতার
সংজ্ঞা থেকে সহজেই দেখা যায়, $\Th{Q}$-কে প্রসারিত করে এমন (অথবা যার মধ্যে
$\Th{Q}$-কে ব্যাখ্যা করা যায় এমন) যেকোনো তত্ত্ব গণনসাধ্য অপেক্ষকগুলিকে
উপস্থাপন করতে পারে। বিপরীত দিকে, যে সঙ্গতিপূর্ণ !!{derivation} পদ্ধতিতে
!!{derivation}-এর ধারণাটি গণনসাধ্য, সেখানে প্রত্যেক উপস্থাপনযোগ্য অপেক্ষক
গণনসাধ্য। তাই, যেমন, গণনসাধ্য অপেক্ষকের সেটকে পেয়ানো পাটিগণিত, এমনকি
জার্মেলো--ফ্রেঙ্কেল সেটতত্ত্বে উপস্থাপনযোগ্য অপেক্ষকের সেট হিসেবেও চিহ্নিত
করা যায়। গ্যোডেল লক্ষ করেছিলেন, এটি কিছুটা বিস্ময়কর। প্রমাণযোগ্যতার ক্ষেত্রে
কোন তত্ত্ব বিবেচনা করছি, প্রশ্নগুলি তার প্রতি খুব সংবেদনশীল—এ কথা আমরা দেখব;
মোটামুটিভাবে, স্বতঃসিদ্ধ যত শক্তিশালী, তত বেশি প্রমাণ করা যায়। কিন্তু
সঙ্গতিপূর্ণ স্বতঃসিদ্ধায়িত তত্ত্বের এক বিস্তৃত শ্রেণি জুড়ে উপস্থাপনযোগ্য
অপেক্ষক ঠিক গণনসাধ্য অপেক্ষকগুলিই; শক্তিশালী তত্ত্ব বেশি অপেক্ষক উপস্থাপন
করে না।
% BN-SRC-225 TARGET-CORRECTION-BEGIN
\par\smallskip\noindent\textnormal{\small\textbf{উৎস-সংশোধন (BN-SRC-225):} হিমায়িত সাধারণীকরণে কেবল নিষ্পাদন গণনসাধ্য এবং তত্ত্ব স্বতঃসিদ্ধায়িত হওয়ার শর্ত দেওয়া হয়েছিল। কিন্তু অসঙ্গতিপূর্ণ তত্ত্বে যেকোনো সূত্র ও তার অস্বীকার উভয়ই প্রমাণযোগ্য, ফলে উপস্থাপনযোগ্যতার সংজ্ঞা অগণনসাধ্য অপেক্ষকের জন্যও তুচ্ছভাবে পূর্ণ হয়। বাংলা পাঠে উভয় সাধারণ দাবিতে আবশ্যক সঙ্গতির শর্ত যোগ করা হয়েছে।} % BN-SRC-225 TARGET-CORRECTION-END
\end{explain}

\end{document}
